\input{../tex-inputs/latex-leseansicht-vorspann}

\section[Arthur Schnitzler an Gustav Schwarzkopf, 6. 9. 1902]{L04412 Arthur Schnitzler an Gustav Schwarzkopf, 6. 9. 1902}
\nopagebreak\mylabel{L04412v}
\rehead{ }\normalsize\beginnumbering\briefempfaengerindex{Schwarzkopf, Gustav@\textsc{Schwarzkopf, Gustav}!zzzSchnitzler, Arthur@\emph{von Arthur Schnitzler}!1902-09-061@{6. 9. 1902}|(be}
\toendnotes[C]{\smallbreak\pagebreak[2]}
\correspDesc{Versand  durch Arthur Schnitzler am 6. 9. 1902 in Windischgarsten
\newline{}Übermittlung  am 6. 9. 1902 in Windischgarsten
\newline{}Zustellung  am 7. 9. 1902 in Wien 1/1 1
\newline{}Erhalt  durch Gustav Schwarzkopf am 7. 9. 1902 in Tiefer Graben 23}\toendnotes[C]{\smallbreak}
\Standort{DLA, A:Schnitzler, HS.1985.1.1897.}
\physDesc{Bildpostkarte, 62 Zeichen
\newline{}Handschrift: Bleistift, deutsche Kurrent
\newline{}Versand: 1) Stempel: »\nobreak{}\oindex{Windischgarsten@\textbf{Windischgarsten}|pwk}Windisch\textcolor{gray}{garsten}, 6. 9. 0{[}2{]}\nobreak{}«.   2) Stempel: »\nobreak{}\oindex{I., Innere Stadt@\textbf{I., Innere Stadt}|pwk}Wien 1/1 \textcolor{gray}{1}, 7. \textcolor{gray}{9. 1904}, 9–10½V., Bestellt\nobreak{}«. }\toendnotes[C]{\smallbreak}\pstart{}{\pb}Herrn \textsc{Gustav Schwarzkopf}\pend{}\pstart{}Wien\oindex{Wien@\textbf{Wien}|pw}\pend{}\pstart{}\textsc{I. Tiefer Graben 23}\oindex{Wien@\textbf{Wien}!I., Innere Stadt@\textbf{I., Innere Stadt}!Tiefer Graben 23@\textbf{Tiefer Graben 23}|pw}\pend{}{\bigskip}
\pstart
           \centering{}{\pb}\textcolor{gray}{\textbf{\label{K_L04412-1v}\edtext{Hinterstoder}{\lemma{\textnormal{\emph{Hinterstoder}}}\Cendnote{\textnormal{Obzwar nur durch den Poststempel
                        datiert, lässt sich diese Karte durch den Abgleich mit dem \emph{Tagebuch}\pwindex{Schnitzler, Arthur 15. 5. 1862 Wien – 21. 10. 1931 ebd.@\textsc{Schnitzler, Arthur} (15. 5. 1862 Wien – 21. 10. 1931 ebd.)!Tagebuch@\strich\emph{Tagebuch}|pwk} vom 6. 9. 1902 verlässlich datieren.}}}\label{K_L04412-1}
                     (Ob.-Öst)\oindex{Hinterstoder@\textbf{Hinterstoder}|pw} mit Villa Herzog von
                     Würtemberg\oindex{Villa Württemberg@\textbf{Villa Württemberg}|pw}.}}\pend
           \vspace{1em}
\pstart
           {\pb}Herzlich Gruß.{\\[\baselineskip]}\spacefill\mbox{A.}\pend
           \leftskip=0em{}\selectlanguage{ngerman}\endnumbering\briefempfaengerindex{Schwarzkopf, Gustav@\textsc{Schwarzkopf, Gustav}!zzzSchnitzler, Arthur@\emph{von Arthur Schnitzler}!1902-09-061@{6. 9. 1902}|)be}\mylabel{L04412h}
\begin{anhang}
\end{anhang}\newcommand{\dateiname}{L04412}\newcommand{\titel}{Arthur Schnitzler an Gustav Schwarzkopf, 6. 9. 1902}\newcommand{\editorInnen}{Herausgegeben von Jahnke, SelmaMüller, Martin Anton}\input{../tex-inputs/latex-leseansicht-abspann}
